\documentclass{infi-ex}

\exnumber{2}
\extopic{המספרים הממשיים}

\begin{document}
\makeextitle

\begin{instructions}
יש להוכיח את הטענות הבאות בעזרת האקסיומות של הממשיים בלבד (בסוף התרגיל מצורפת הטבלה מהכיתה).

\textbf{הערה:} מותר להשתמש בטענות שהוכחו בהרצאה או בתרגול בלי להוכיח אותן, אבל יש לציין איפה ראיתם את ההוכחה.
\end{instructions}

\begin{question}[קיום איבר נייטרלי יחיד לחיבור]
\[ \forall x.\forall y.\ \bigl((\forall z.\ x+z=z\land y+z=y)\to x=y\bigr) \]
\end{question}

\begin{question}[האיבר ההופכי לכפל של מספר חיובי הוא גם חיובי]
\[ \forall x.\ 0<x\to0<x^{-1} \]
\end{question}

\begin{question}[שלילת השלילה]
\[ \forall x.\ -(-x)=x \]
\end{question}

\begin{question}[היפוך הסדר של ההופכי לחיבור]
\[ \forall x.\forall y.\ x\le y\to -y\le -x \]
\end{question}

\begin{question}[כלל הצמצום בחיבור]
\[ \forall x.\forall y.\forall z.\ x+z=y+z\to x=y \]
\end{question}

\begin{question}[כפל במינוס אחד]
\[ \forall x.\ (-1)\cdot x=-x \]
\end{question}

\begin{question}[טרנזיטיביות של $<$]
\[ \forall x.\forall y.\forall z.\ (x<y\land y<z)\to x<z \]
\end{question}

% axioms.tex — טבלת אקסיומות הממשיים (מהכיתה). לצירוף לתרגיל: % axioms.tex — טבלת אקסיומות הממשיים (מהכיתה). לצירוף לתרגיל: % axioms.tex — טבלת אקסיומות הממשיים (מהכיתה). לצירוף לתרגיל: \input{axioms}
\newpage
\begingroup
\renewcommand\arraystretch{1.2}
\newcommand\axsec[1]{\par\addvspace{5pt}{\sffamily\bfseries\color{exaccent}#1}\par\nopagebreak\vspace{2pt}}
\newenvironment{axtable}{%
  \noindent\small
  \rowcolors{2}{white}{exlight}%
  \begin{tabular}{>{\raggedright}p{0.23\linewidth}>{\raggedright}p{0.45\linewidth}>{\raggedright\arraybackslash}p{0.24\linewidth}}
  \rowcolor{exaccent}
  \textcolor{white}{\sffamily\bfseries שם} &
  \textcolor{white}{\sffamily\bfseries אקסיומה} &
  \textcolor{white}{\sffamily\bfseries הערות} \\}
  {\end{tabular}\par}

\begin{center}
  {\sffamily\Large\bfseries\color{exaccent}נספח: אקסיומות המספרים הממשיים}
\end{center}

\axsec{סימנים}
\begin{itemize}
  \item לוגיים: $\land,\ \lor,\ \lnot,\ \to,\ \leftrightarrow,\ \forall,\ \exists,\ =$
  \item קבוצות: $\in,\ \subseteq,\ \emptyset$
  \item ממשיים: $+,\ \cdot,\ \le,\ 0,\ 1$, משתנים $x,y,z,\dots$
\end{itemize}

\axsec{סימונים וקיצורים}
\begin{tabular}{p{0.3\linewidth}p{0.3\linewidth}p{0.3\linewidth}}
$x<y$ קיצור של $x\le y\land x\ne y$ &
$x>y$ קיצור של $y<x$ &
$x\ge y$ קיצור של $y\le x$ \\
$x-y$ קיצור של $x+(-y)$ &
$\frac xy$ קיצור של $x\cdot y^{-1}$ &
$xy$ קיצור של $x\cdot y$
\end{tabular}

\axsec{אקסיומות החיבור}
\begin{axtable}
איבר נייטרלי לחיבור & $\forall y.\ y+0=y$ & \\
איבר הופכי לחיבור & $\forall x.\exists y.\ x+y=0$ & האיבר $y$ מסומן ב־\mbox{$-x$} \\
אסוציאטיביות החיבור & $\forall x.\forall y.\forall z.\ (x+y)+z=x+(y+z)$ & \\
קומוטטיביות החיבור & $\forall x.\forall y.\ x+y=y+x$ & \\
\end{axtable}

\axsec{אקסיומות הכפל}
\begin{axtable}
איבר נייטרלי לכפל & $1\ne0\land\forall y.\ 1\cdot y=y$ & \\
איבר הופכי לכפל & $\forall x.\ (x\ne0\to\exists y.\ x\cdot y=1)$ & האיבר $y$ מסומן ב־\mbox{$x^{-1}$} \\
אסוציאטיביות הכפל & $\forall x.\forall y.\forall z.\ x\cdot(y\cdot z)=(x\cdot y)\cdot z$ & \\
קומוטטיביות הכפל & $\forall x.\forall y.\ x\cdot y=y\cdot x$ & \\
\end{axtable}

\axsec{האקסיומה המקשרת בין כפל וחיבור}
\begin{axtable}
דיסטריביוטיביות & $\forall x.\forall y.\forall z.\ x\cdot(y+z)=(x\cdot y)+(x\cdot z)$ & \\
\end{axtable}

\axsec{אקסיומות יחס הסדר}
\begin{axtable}
רפלקסיביות & $\forall x.\ x\le x$ & \\
אנטי־סימטריות & $\forall x.\forall y.\ (x\le y\land y\le x)\to x=y$ & \\
טרנזיטיביות & $\forall x.\forall y.\forall z.\ (x\le y\land y\le z)\to x\le z$ & \\
לינאריות & $\forall x.\forall y.\ x\le y\lor y\le x$ & \\
\end{axtable}

\axsec{אקסיומות המקשרות בין הפעולות ויחס הסדר}
\begin{axtable}
מונוטוניות החיבור & $\forall x.\forall y.\forall z.\ x\le y\to x+z\le y+z$ & \\
מונוטוניות הכפל & $\forall x.\forall y.\ (x\ge0\land y\ge0)\to x\cdot y\ge0$ & \\
\end{axtable}

\axsec{אקסיומת השלמות}
\begin{axtable}
שלמות &
$\forall X.\forall Y.\ \bigl(X\ne\emptyset\land Y\ne\emptyset\land\forall x\in X.\forall y\in Y.\ x\le y\bigr)$
$\to\bigl(\exists c.\forall x\in X.\forall y\in Y.\ x\le c\le y\bigr)$ &
האקסיומה החשובה ביותר \\
\end{axtable}
\endgroup

\newpage
\begingroup
\renewcommand\arraystretch{1.2}
\newcommand\axsec[1]{\par\addvspace{5pt}{\sffamily\bfseries\color{exaccent}#1}\par\nopagebreak\vspace{2pt}}
\newenvironment{axtable}{%
  \noindent\small
  \rowcolors{2}{white}{exlight}%
  \begin{tabular}{>{\raggedright}p{0.23\linewidth}>{\raggedright}p{0.45\linewidth}>{\raggedright\arraybackslash}p{0.24\linewidth}}
  \rowcolor{exaccent}
  \textcolor{white}{\sffamily\bfseries שם} &
  \textcolor{white}{\sffamily\bfseries אקסיומה} &
  \textcolor{white}{\sffamily\bfseries הערות} \\}
  {\end{tabular}\par}

\begin{center}
  {\sffamily\Large\bfseries\color{exaccent}נספח: אקסיומות המספרים הממשיים}
\end{center}

\axsec{סימנים}
\begin{itemize}
  \item לוגיים: $\land,\ \lor,\ \lnot,\ \to,\ \leftrightarrow,\ \forall,\ \exists,\ =$
  \item קבוצות: $\in,\ \subseteq,\ \emptyset$
  \item ממשיים: $+,\ \cdot,\ \le,\ 0,\ 1$, משתנים $x,y,z,\dots$
\end{itemize}

\axsec{סימונים וקיצורים}
\begin{tabular}{p{0.3\linewidth}p{0.3\linewidth}p{0.3\linewidth}}
$x<y$ קיצור של $x\le y\land x\ne y$ &
$x>y$ קיצור של $y<x$ &
$x\ge y$ קיצור של $y\le x$ \\
$x-y$ קיצור של $x+(-y)$ &
$\frac xy$ קיצור של $x\cdot y^{-1}$ &
$xy$ קיצור של $x\cdot y$
\end{tabular}

\axsec{אקסיומות החיבור}
\begin{axtable}
איבר נייטרלי לחיבור & $\forall y.\ y+0=y$ & \\
איבר הופכי לחיבור & $\forall x.\exists y.\ x+y=0$ & האיבר $y$ מסומן ב־\mbox{$-x$} \\
אסוציאטיביות החיבור & $\forall x.\forall y.\forall z.\ (x+y)+z=x+(y+z)$ & \\
קומוטטיביות החיבור & $\forall x.\forall y.\ x+y=y+x$ & \\
\end{axtable}

\axsec{אקסיומות הכפל}
\begin{axtable}
איבר נייטרלי לכפל & $1\ne0\land\forall y.\ 1\cdot y=y$ & \\
איבר הופכי לכפל & $\forall x.\ (x\ne0\to\exists y.\ x\cdot y=1)$ & האיבר $y$ מסומן ב־\mbox{$x^{-1}$} \\
אסוציאטיביות הכפל & $\forall x.\forall y.\forall z.\ x\cdot(y\cdot z)=(x\cdot y)\cdot z$ & \\
קומוטטיביות הכפל & $\forall x.\forall y.\ x\cdot y=y\cdot x$ & \\
\end{axtable}

\axsec{האקסיומה המקשרת בין כפל וחיבור}
\begin{axtable}
דיסטריביוטיביות & $\forall x.\forall y.\forall z.\ x\cdot(y+z)=(x\cdot y)+(x\cdot z)$ & \\
\end{axtable}

\axsec{אקסיומות יחס הסדר}
\begin{axtable}
רפלקסיביות & $\forall x.\ x\le x$ & \\
אנטי־סימטריות & $\forall x.\forall y.\ (x\le y\land y\le x)\to x=y$ & \\
טרנזיטיביות & $\forall x.\forall y.\forall z.\ (x\le y\land y\le z)\to x\le z$ & \\
לינאריות & $\forall x.\forall y.\ x\le y\lor y\le x$ & \\
\end{axtable}

\axsec{אקסיומות המקשרות בין הפעולות ויחס הסדר}
\begin{axtable}
מונוטוניות החיבור & $\forall x.\forall y.\forall z.\ x\le y\to x+z\le y+z$ & \\
מונוטוניות הכפל & $\forall x.\forall y.\ (x\ge0\land y\ge0)\to x\cdot y\ge0$ & \\
\end{axtable}

\axsec{אקסיומת השלמות}
\begin{axtable}
שלמות &
$\forall X.\forall Y.\ \bigl(X\ne\emptyset\land Y\ne\emptyset\land\forall x\in X.\forall y\in Y.\ x\le y\bigr)$
$\to\bigl(\exists c.\forall x\in X.\forall y\in Y.\ x\le c\le y\bigr)$ &
האקסיומה החשובה ביותר \\
\end{axtable}
\endgroup

\newpage
\begingroup
\renewcommand\arraystretch{1.2}
\newcommand\axsec[1]{\par\addvspace{5pt}{\sffamily\bfseries\color{exaccent}#1}\par\nopagebreak\vspace{2pt}}
\newenvironment{axtable}{%
  \noindent\small
  \rowcolors{2}{white}{exlight}%
  \begin{tabular}{>{\raggedright}p{0.23\linewidth}>{\raggedright}p{0.45\linewidth}>{\raggedright\arraybackslash}p{0.24\linewidth}}
  \rowcolor{exaccent}
  \textcolor{white}{\sffamily\bfseries שם} &
  \textcolor{white}{\sffamily\bfseries אקסיומה} &
  \textcolor{white}{\sffamily\bfseries הערות} \\}
  {\end{tabular}\par}

\begin{center}
  {\sffamily\Large\bfseries\color{exaccent}נספח: אקסיומות המספרים הממשיים}
\end{center}

\axsec{סימנים}
\begin{itemize}
  \item לוגיים: $\land,\ \lor,\ \lnot,\ \to,\ \leftrightarrow,\ \forall,\ \exists,\ =$
  \item קבוצות: $\in,\ \subseteq,\ \emptyset$
  \item ממשיים: $+,\ \cdot,\ \le,\ 0,\ 1$, משתנים $x,y,z,\dots$
\end{itemize}

\axsec{סימונים וקיצורים}
\begin{tabular}{p{0.3\linewidth}p{0.3\linewidth}p{0.3\linewidth}}
$x<y$ קיצור של $x\le y\land x\ne y$ &
$x>y$ קיצור של $y<x$ &
$x\ge y$ קיצור של $y\le x$ \\
$x-y$ קיצור של $x+(-y)$ &
$\frac xy$ קיצור של $x\cdot y^{-1}$ &
$xy$ קיצור של $x\cdot y$
\end{tabular}

\axsec{אקסיומות החיבור}
\begin{axtable}
איבר נייטרלי לחיבור & $\forall y.\ y+0=y$ & \\
איבר הופכי לחיבור & $\forall x.\exists y.\ x+y=0$ & האיבר $y$ מסומן ב־\mbox{$-x$} \\
אסוציאטיביות החיבור & $\forall x.\forall y.\forall z.\ (x+y)+z=x+(y+z)$ & \\
קומוטטיביות החיבור & $\forall x.\forall y.\ x+y=y+x$ & \\
\end{axtable}

\axsec{אקסיומות הכפל}
\begin{axtable}
איבר נייטרלי לכפל & $1\ne0\land\forall y.\ 1\cdot y=y$ & \\
איבר הופכי לכפל & $\forall x.\ (x\ne0\to\exists y.\ x\cdot y=1)$ & האיבר $y$ מסומן ב־\mbox{$x^{-1}$} \\
אסוציאטיביות הכפל & $\forall x.\forall y.\forall z.\ x\cdot(y\cdot z)=(x\cdot y)\cdot z$ & \\
קומוטטיביות הכפל & $\forall x.\forall y.\ x\cdot y=y\cdot x$ & \\
\end{axtable}

\axsec{האקסיומה המקשרת בין כפל וחיבור}
\begin{axtable}
דיסטריביוטיביות & $\forall x.\forall y.\forall z.\ x\cdot(y+z)=(x\cdot y)+(x\cdot z)$ & \\
\end{axtable}

\axsec{אקסיומות יחס הסדר}
\begin{axtable}
רפלקסיביות & $\forall x.\ x\le x$ & \\
אנטי־סימטריות & $\forall x.\forall y.\ (x\le y\land y\le x)\to x=y$ & \\
טרנזיטיביות & $\forall x.\forall y.\forall z.\ (x\le y\land y\le z)\to x\le z$ & \\
לינאריות & $\forall x.\forall y.\ x\le y\lor y\le x$ & \\
\end{axtable}

\axsec{אקסיומות המקשרות בין הפעולות ויחס הסדר}
\begin{axtable}
מונוטוניות החיבור & $\forall x.\forall y.\forall z.\ x\le y\to x+z\le y+z$ & \\
מונוטוניות הכפל & $\forall x.\forall y.\ (x\ge0\land y\ge0)\to x\cdot y\ge0$ & \\
\end{axtable}

\axsec{אקסיומת השלמות}
\begin{axtable}
שלמות &
$\forall X.\forall Y.\ \bigl(X\ne\emptyset\land Y\ne\emptyset\land\forall x\in X.\forall y\in Y.\ x\le y\bigr)$
$\to\bigl(\exists c.\forall x\in X.\forall y\in Y.\ x\le c\le y\bigr)$ &
האקסיומה החשובה ביותר \\
\end{axtable}
\endgroup


\end{document}
